% Complete theoremhood and decision treatment for the existing X/strong-U sheet.
\section{Linear Temporal Completeness and Decision}
\label{sec:p08-linear-temporal-completeness}

This chapter proves the adequacy of
\hyperref[sheet:hilbert-ltl]{the registered Hilbert LTL calculus}
on discrete infinite words, then constructs an accepting-run decision
procedure. Its until operator permits a witness at the current position.
The distinction from a strict future operator, and from consequence
using infinitely many premises, matters to both proofs.

The historical reference for finite-model threading is Gabbay,
Pnueli, Shelah and Stavi's published 1980 paper. That paper's
strict-future language and DUX calculus are kept distinct below.
The automaton construction uses the local-type and pending-batch
method in Appendix B of Vardi and Wolper's author preprint dated
15 August 1985, associated with their 1986 proceedings publication.
Sistla and Clarke's 1985 journal article supplies the complexity
comparison. The proofs here fill the stated construction steps for
the exact registered language; a reformulation is not presented as
a verbatim reproduction of another calculus or source version
\cite{GabbayPnueliShelahStavi1980,VardiWolper1986,SistlaClarke1985}.

\subsection{The exact language and premise boundary}
\label{sec:p08-ltl-completeness-language}

\begin{hbdefinition}{def:p08-ltl-completeness-language}{The fixed infinite-word problem}
Use the atom set \(\mathsf{Prop}\) and the grammar
\[
 A::=p\mid\bottom\mid\neg A\mid(A\land A)
       \mid\mathsf X A\mid A\mathbin{\mathsf U}A.
\]
Boolean abbreviations are those in
\cref{sec:modal-temporal-epistemic-02}; in particular
\(\mathsf F A=\topv\mathbin{\mathsf U}A\) and
\(\mathsf G A=\neg\mathsf F\neg A\).
A word is \(\mathcal M=(\mathbb N,V)\) with
\(V:\mathbb N\to\mathcal P(\mathsf{Prop})\). Next means truth at
\(n+1\), and
\[
 \mathcal M,n\entails A\mathbin{\mathsf U}B
 \quad\Longleftrightarrow\quad
 \exists j\geq n\,
 \bigl(\mathcal M,j\entails B\ \land\
       \forall k\,(n\leq k<j\Rightarrow\mathcal M,k\entails A)\bigr).
\]
The Boolean clauses are classical. Write \(\vdash_L A\) for the
existing finite no-premise Hilbert derivation
\(\derives_{\mathsf{LTL}}A\). Its temporal schemas are X-distribution,
X-negation and U-unfolding; its rules are MP, theorem-only XNec
and theorem-only UInd, exactly as registered. No rule here allows
necessitating or applying UInd to an open assumption.
\end{hbdefinition}

\begin{hbtheorem}{thm:p08-ltl-strict-bridge}{proved=annex:p08.ltl.strict-bridge}{The strict and non-strict conventions on natural-number time}
Let \(A\mathbin{\mathsf U_s}B\) mean that a witness \(j>n\)
satisfies \(B\), with \(A\) required only at \(n<k<j\).
This is the strict convention used in the 1980 paper. On these words,
\[
 A\mathbin{\mathsf U_s}B\ \leftrightarrow\
 \mathsf X(A\mathbin{\mathsf U}B),\qquad
 A\mathbin{\mathsf U}B\ \leftrightarrow\
 B\lor(A\land(A\mathbin{\mathsf U_s}B)).
\]
The displayed biconditionals assert pointwise truth equivalence
of the explicitly compared expressions. Replacing operators
recursively gives truth-preserving translations of nested formulas.
Strict eventuality and invariance correspond to
\(\mathsf X\mathsf F A\) and \(\mathsf X\mathsf G A\).
In particular, \(\bottom\mathbin{\mathsf U_s}B\) means next \(B\),
whereas \(\bottom\mathbin{\mathsf U}B\) means current \(B\).
\end{hbtheorem}

\subsection{Soundness and the premise boundary}
\label{sec:p08-ltl-soundness-boundary}

\begin{hbtheorem}{thm:p08-ltl-soundness}{proved=annex:p08.ltl.soundness}{Soundness of every registered schema and rule}
If \(\vdash_L A\), then \(\mathcal M,n\entails A\) for every
word and every \(n\). Moreover a formula is valid at all positions
of all words exactly when it is valid at position zero of all words.
\end{hbtheorem}

\begin{hbdefinition}{def:p08-ltl-local-consequence}{Local premise consequence and its finitary proof relation}
Write \(\Gamma\models_\omega A\) when, at every word and position,
truth of every member of \(\Gamma\) implies truth of \(A\) there.
Define \(\Gamma\vdash_L^{\rm fin}A\) to mean that some finite
\(\Gamma_0\subseteq\Gamma\) satisfies
\(\vdash_L(\bigwedge\Gamma_0)\to A\), with empty conjunction
\(\topv\). Premises are local to the evaluation position, not
assumed at every later position. This relation does not change
the registered theorem-only rules.
\end{hbdefinition}

\begin{hbtheorem}{thm:p08-ltl-noncompactness}{proved=annex:p08.ltl.noncompactness}{The exact arbitrary-premise failure}
For an atom \(p\), put
\(\Gamma=\{\mathsf X^n p:n\in\mathbb N\}\), with
\(\mathsf X^0p=p\). Then \(\Gamma\models_\omega\mathsf Gp\),
but no finite subset entails \(\mathsf Gp\). Equivalently,
\(\Gamma\cup\{\neg\mathsf Gp\}\) is unsatisfiable while every
finite subset is satisfiable. Thus sound finitary derivations cannot
be strongly complete for this arbitrary-premise consequence relation.
\end{hbtheorem}

\subsection{Finite consistent types and their successors}
\label{sec:p08-ltl-consistent-types}

\begin{hblemma}{thm:p08-ltl-next-boolean}{proved=annex:p08.ltl.next-boolean}{Derived Boolean laws for next}
The registered calculus proves that \(\mathsf X\) preserves truth,
falsity, negation, finite conjunctions and finite disjunctions, and
preserves every provable implication. In particular, from
\(\vdash_L C\to D\) one obtains
\(\vdash_L\mathsf XC\to\mathsf XD\); no premises are open
in this derivation.
\end{hblemma}

\begin{hbdefinition}{def:p08-ltl-consistent-type}{Finite truth records and characteristic formulas}
For a formula \(\varphi\), expand abbreviations and let
\(C=\operatorname{Sub}(\varphi)\), its finite set of syntactic
subformulas. Let \(n\) be the length of that expanded formula,
counting its written symbol occurrences, and put \(m=|C|\).
Atom identifiers are effectively encoded; counting their full
written lengths can only increase this input-size bound.
A truth record is a function
\(t:C\to\{0,1\}\). Its characteristic formula is
\[
 \chi_t=\bigwedge_{D\in C}D^{t(D)},\qquad
 D^1=D,\quad D^0=\neg D,
\]
in a fixed order. Call \(t\) \emph{consistent} when
\(\nvdash_L\neg\chi_t\), and let \(Q_C\) be the set of such
records. Define \(tR_Cs\) exactly when
\(\nvdash_L\neg(\chi_t\land\mathsf X\chi_s)\).
This finite graph is used for the existence proof. Testing its
definition is not the algorithm later used for deciding formulas.
\end{hbdefinition}

\begin{hblemma}{thm:p08-ltl-type-partition}{proved=annex:p08.ltl.type-partition}{A finite exhaustive partition}
The calculus proves \(\bigvee_{t\in Q_C}\chi_t\), and different
characteristic formulas are pairwise incompatible. Each consistent
record satisfies the Boolean truth conditions on \(C\). If
\(\nvdash_L\neg\varphi\), some \(t\in Q_C\) has
\(t(\varphi)=1\). Empty disjunction means \(\bottom\).
\end{hblemma}

\begin{hblemma}{thm:p08-ltl-successor}{proved=annex:p08.ltl.successor}{Successor existence and all one-step obligations}
Every \(t\in Q_C\) has an \(R_C\)-successor. On every edge
\(tR_Cs\),
\[
\begin{aligned}
 t(\mathsf XD)&=s(D)&& (\mathsf XD\in C),\\
 t(A\mathbin{\mathsf U}B)&=
 \max\{t(B),\min(t(A),s(A\mathbin{\mathsf U}B))\}
 &&(A\mathbin{\mathsf U}B\in C).
\end{aligned}
\]
An until formula's extra \(\mathsf X\)-obligation need not be
inserted into \(C\); the unfolding axiom and the edge's consistency
already enforce it.
\end{hblemma}

\begin{hblemma}{thm:p08-ltl-eventuality-reach}{proved=annex:p08.ltl.eventuality-reach}{UInd forces a finite witness path}
If \(t(A\mathbin{\mathsf U}B)=1\), there is a finite path
\(t=t_0R_Ct_1R_C\cdots R_Ct_k\) with
\(t_k(B)=1\) and \(t_i(A)=1\) for every \(i<k\).
Length zero is permitted when the right operand already holds.
\end{hblemma}

\subsection{Realizing all eventualities and proving completeness}
\label{sec:p08-ltl-fair-truth}

\begin{hbdefinition}{def:p08-ltl-fair-sequence}{Local coherence and eventuality realization}
A sequence \(t_0,t_1,\ldots\) of records on \(C\) is locally
coherent when each record satisfies the Boolean conditions and
every consecutive pair satisfies the two displayed one-step clauses.
It realizes eventualities when, whenever
\(t_i(A\mathbin{\mathsf U}B)=1\), some \(j\geq i\) has
\(t_j(B)=1\). The one-step clauses then force \(A\) at every
preceding point of a least such witness interval.
\end{hbdefinition}

\begin{hblemma}{thm:p08-ltl-fair-path}{proved=annex:p08.ltl.fair-path}{A single infinite path realizes every pending until}
From every consistent record there is an infinite \(R_C\)-path
realizing all eventualities. The construction uses a fixed finite
ordering of the until subformulas and finite witness paths, with
an extra successor step in every scheduling round.
\end{hblemma}

\begin{hblemma}{thm:p08-ltl-fair-truth}{proved=annex:p08.ltl.fair-truth}{The full truth lemma for a coherent realizing sequence}
For any locally coherent sequence realizing eventualities, define
\(V(i)\) on atoms occurring in \(\varphi\) by
\(p\in V(i)\) exactly when \(t_i(p)=1\), and set unused atoms
false. Then for all \(D\in C\) and all \(i\),
\[
 \mathcal M,i\entails D\quad\Longleftrightarrow\quad t_i(D)=1.
\]
Both directions for until are part of the claim; local unfolding
alone is insufficient.
\end{hblemma}

\begin{hbtheorem}{thm:p08-linear-temporal-completeness}{proved=annex:p08.ltl.completeness}{Weak and finite-premise completeness for the registered calculus}
For every formula \(A\),
\[
 \vdash_L A\quad\Longleftrightarrow\quad
 \mathcal M,n\entails A\text{ for every word and position}.
\]
For every finite \(\Gamma\),
\(\Gamma\models_\omega A\) exactly when
\(\vdash_L(\bigwedge\Gamma)\to A\).
The arbitrary-premise noncompactness theorem remains applicable.
\end{hbtheorem}

\begin{hbexample}{def:p08-ltl-unfulfilled-loop}{Why a locally consistent loop need not be a model}
For \(\varphi=\mathsf Fp\), a local record can make
\(p\) false and \(\mathsf Fp\) true forever while satisfying
\(\mathsf Fp\leftrightarrow p\lor\mathsf X\mathsf Fp\)
at every step. The resulting loop has no witness for \(p\) and
therefore fails eventuality realization. By contrast, a record
sequence in which \(p\) first becomes true at position three
realizes \(\mathsf Fp\) at position zero. Recurrence
\(\mathsf G\mathsf Fp\) requires further witnesses indefinitely.
\end{hbexample}

\subsection{An effective local automaton and its pending batches}
\label{sec:p08-ltl-automaton-construction}

\begin{hbdefinition}{def:p08-ltl-buchi-convention}{Infinite-word automata with permitted starts}
A nondeterministic B\"uchi automaton is
\(\mathcal A=(\Sigma,Q,I,\Delta,F)\), with finite alphabet
and state set, \(I,F\subseteq Q\), and
\(\Delta:Q\times\Sigma\to\mathcal P(Q)\).
On \(a_0a_1\cdots\), a run is \(q_0q_1\cdots\) with
\(q_0\in I\) and \(q_{i+1}\in\Delta(q_i,a_i)\).
Acceptance requires \(q_i\in F\) for infinitely many indices.
The permitted-start set matches the source construction; its
conversion to the handbook's single-start convention is proved below.
One accepting visit, including the initial visit, is insufficient.
\end{hbdefinition}

\begin{hbdefinition}{def:p08-ltl-local-automaton}{The effective closure-type construction}
For input \(\varphi\), let \(C=\operatorname{Sub}(\varphi)\)
as above, \(P\) its finite atom set, and \(\mathcal U\) its
until subformulas. Write \(m=|C|\), \(r=|\mathcal U|\).
Let \(Q_b\) contain \emph{all} records \(t:C\to\{0,1\}\)
satisfying only the Boolean truth conditions:
\(t(\bottom)=0\), \(t(\neg D)=1-t(D)\), and
\(t(D\land E)=\min\{t(D),t(E)\}\) whenever the displayed
compound belongs to \(C\). Define \(tEs\)
by the X and until equations in
\cref{thm:p08-ltl-successor}. This test uses finite truth bits,
not Hilbert consistency. Put
\[
 D_t=\{A\mathbin{\mathsf U}B\in\mathcal U:
                 t(A\mathbin{\mathsf U}B)=1, t(B)=0\},
 \qquad H_t=\{A\mathbin{\mathsf U}B\in\mathcal U:t(B)=1\}.
\]
An obligation state is any \(O\subseteq\mathcal U\).
Its allowed next sets satisfy
\[
 O'\supseteq D_t\quad\text{if }O=\varnothing,
 \qquad O'\supseteq O\setminus H_t\quad\text{if }O\ne\varnothing.
\]
Arbitrary extra obligations are permitted, as in the source;
the accepting choices need not add them.

The automaton \(\mathcal A_\varphi\) has alphabet
\(2^P\), states \(Q_b\times2^{\mathcal U}\), starts
\(\{t:t(\varphi)=1\}\times\{\varnothing\}\), and accepting
set \(Q_b\times\{\varnothing\}\). A transition on letter
\(a\subseteq P\) from \((t,O)\) to \((s,O')\) requires
\(a=\{p\in P:t(p)=1\}\), \(tEs\), and the displayed
obligation-set condition. Thus the current type describes the
letter consumed by its outgoing transition.
\cite[preprint Appendix B, pp.~2--3]{VardiWolper1985LTLPreprint}
\end{hbdefinition}

These truth-bit records encode the source's sets of subformulas
and their signed negations: include \(D\) when its bit is one,
otherwise include \(\neg D\), and identify double negations.
The Boolean conditions make this a bijection with the source's
complete consistent Boolean types. The source first reads whole
closure types, then projects a letter to its true atoms. The
displayed transition performs that projection directly.

\begin{hblemma}{thm:p08-ltl-batch-realization}{proved=annex:p08.ltl.batch-realization}{Exactly what recurring empty obligation sets guarantee}
For a locally coherent type sequence, an allowed obligation
sequence starting at \(\varnothing\) and visiting
\(\varnothing\) infinitely often exists exactly when the type
sequence realizes every eventuality. This includes the case
\(\mathcal U=\varnothing\).
\end{hblemma}

\begin{hbtheorem}{thm:p08-ltl-automaton-language}{proved=annex:p08.ltl.automaton-language}{Both language inclusions and the finite size bound}
The language of \(\mathcal A_\varphi\) is exactly the words
satisfying \(\varphi\) at zero. It has at most
\(2^{m+r}\leq2^{2n}\) states, where \(n\) is the length of
the expanded input formula. This Boolean-coordinate bound implies
the source's looser bound \(2^{3n}\). The construction is
effective, including its initial, transition and accepting sets.
No determinization or complementation theorem is imported.
\end{hbtheorem}

\subsection{Generalized acceptance and the counter construction}
\label{sec:p08-ltl-generalized-acceptance}

The same finite eventuality requirement can be recorded by several
accepting sets, one for each until, rather than by a set of pending
obligations. The following finite constructions establish the
relation between these two representations; the pending-batch
construction above retains the source's own representation.

\begin{hbdefinition}{def:p08-ltl-generalized-buchi}{Generalized B\"uchi acceptance}
A generalized B\"uchi automaton has the same finite alphabet,
states, permitted starts and transition function as
\cref{def:p08-ltl-buchi-convention}, and a finite list
\(F_1,\ldots,F_h\) of accepting subsets. Its infinite run is
accepting when it visits \emph{each} \(F_j\) infinitely often.
When \(h=0\), every infinite run is accepting.

For the LTL local automaton use \(Q_b\), the starts with
\(t(\varphi)=1\), and the letter-matching edges \(tEs\)
from \cref{def:p08-ltl-local-automaton}. For each
\(U=A\mathbin{\mathsf U}B\in\mathcal U\), put
\[
 F_U=\{t\in Q_b:t(U)=0\text{ or }t(B)=1\}.
\]
Call the resulting generalized automaton \(\mathcal G_\varphi\).
These are finite truth-bit tests; no consistency oracle is used.
\end{hbdefinition}

\begin{hblemma}{thm:p08-ltl-generalized-realization}{proved=annex:p08.ltl.generalized-realization}{Each acceptance set expresses one eventuality}
A locally coherent type sequence realizes every eventuality
exactly when it visits each \(F_U\) infinitely often.
Consequently \(\mathcal G_\varphi\) accepts precisely the
words satisfying \(\varphi\) at zero, and has at most
\(2^m\) states and \(r\) accepting sets. It has the same
language as \(\mathcal A_\varphi\).
\end{hblemma}

\begin{hbdefinition}{def:p08-ltl-buchi-counter}{Tracking the next required accepting set}
For \(h\geq1\), convert a generalized automaton
\((\Sigma,Q,I,\Delta,F_1,\ldots,F_h)\) to an ordinary
B\"uchi automaton with states \(Q\times\{1,\ldots,h\}\),
starts \(I\times\{1\}\), and accepting set
\[
 F^*=\{(q,h):q\in F_h\}.
\]
Its transitions on \(a\) have first coordinate
\(q'\in\Delta(q,a)\), and second coordinate
\[
 j'=\begin{cases}
        j,&q\notin F_j,\\
        j+1,&q\in F_j\text{ and }j<h,\\
        1,&q\in F_j\text{ and }j=h.
      \end{cases}
\]
The index \(j\) is the next accepting set still to be
witnessed; its update tests the current state before consuming
the transition. For \(h=0\), retain the original state and
transition sets and declare all states accepting.
\end{hbdefinition}

\begin{hbtheorem}{thm:p08-ltl-generalized-conversion}{proved=annex:p08.ltl.generalized-conversion}{Language preservation, all edge cases and the state bound}
The counter conversion preserves the infinite-word language
and has \(h|Q|\) states for \(h\geq1\), or \(|Q|\) for
\(h=0\). It handles empty accepting subsets and an empty
start set without an additional assumption. Applied to
\(\mathcal G_\varphi\), it gives an ordinary automaton with
at most \(\max\{1,r\}2^m\) states. A state with counter
\(h\) is accepting only when its first coordinate actually
lies in \(F_h\); merely reaching the last counter value
does not complete a round.
\end{hbtheorem}

\begin{hbtheorem}{thm:p08-ltl-single-start}{proved=annex:p08.ltl.single-start}{The single-start conversion without an empty transition}
For any permitted-start automaton, add a fresh nonaccepting
\(q_*\) with
\[
 \Delta_*(q_*,a)=\bigcup_{q\in I}\Delta(q,a).
\]
Keep every old transition and accepting state, and make
\(\{q_*\}\) the sole start. This preserves the infinite-word
language, including an empty original start set, and adds one
state. No transition targets \(q_*\).
\end{hbtheorem}

\subsection{Accepting cycles and finite-state verification}
\label{sec:p08-ltl-accepting-cycles}

\begin{hbdefinition}{def:p08-ltl-automaton-graph}{Paths, positive returns and eventual periodicity}
The automaton's graph has an edge \(q\to q'\) when some
letter permits that transition. Write \(\to^*\) for a path
of zero or more edges, and \(\to^+\) for one or more edges.
A strongly connected component is an equivalence class for
mutual \(\to^*\)-reachability. It is cyclic when it contains
a nonempty closed walk; a singleton needs a self-loop.
A word \(uv^\omega\), with finite \(u\) and nonempty finite
\(v\), is ultimately periodic.
\end{hbdefinition}

\begin{hbtheorem}{thm:p08-ltl-accepting-cycle}{proved=annex:p08.ltl.accepting-cycle}{Emptiness, a positive return and a lasso witness}
A finite B\"uchi automaton accepts some word exactly when
there are \(i\in I,f\in F\) with
\(i\to^*f\) and \(f\to^+f\). Equivalently, some reachable
cyclic strongly connected component contains an accepting state.
For an automaton with \(N\) states a nonempty language has a
witness \(uv^\omega\) with \(\operatorname{len}(u)\leq N-1\)
and \(1\leq\operatorname{len}(v)\leq N\), where
\(\operatorname{len}\) denotes finite-word length.
\end{hbtheorem}

\begin{hbexample}{def:p08-ltl-positive-return}{An accepting point is not an accepting cycle}
Let \(i\) be the only start and accepting state, and let the
only edges be \(i\to q\) and \(q\to q\), where \(q\)
is nonaccepting. The automaton has an infinite run and starts
in an accepting state, but its language is empty. Counting the
zero-length path \(i\to^*i\) as a return would give the
wrong answer. In the pending-batch construction, a loop with
\(\mathsf Fp\) true and \(p\) always false likewise leaves
its initially empty batch and can never empty it again.
\end{hbexample}

\begin{hbdefinition}{def:p08-ltl-product-checking}{The exact finite-state model-checking problem}
Let \(K=(W,w_0,R,\ell)\) have finite nonempty \(W\),
total relation \(R\), and labels \(\ell:W\to2^P\).
It satisfies \(\varphi\) when every infinite path from
\(w_0\) yields a word satisfying \(\varphi\) at zero.
Construct \(\mathcal A_{\neg\varphi}\). Product vertices
are \((w,t,O)\) whose type's atom bits agree with \(\ell(w)\).
Product edges require \(wRw'\) and the automaton transition
on \(\ell(w)\). Starts use \(w_0\) and the automaton's
starts; accepting vertices have \(O=\varnothing\).
This is universal-path LTL checking, not a CTL formula or a
change to the until operator.
Write \(\lVert K\rVert\geq2\) for the complete encoded
input length of its state, edge and label lists, increased to
two for a smaller representation. It is a size parameter,
not a vector norm or a count of only the reachable states.
\end{hbdefinition}

\begin{hbtheorem}{thm:p08-ltl-product-correctness}{proved=annex:p08.ltl.product-correctness}{Product paths are exactly the counterexamples}
The finite-state structure violates \(\varphi\) exactly when
the product contains a reachable accepting state with a positive
return. Projecting its prefix and loop yields a genuine infinite
counterexample path of \(K\). Totality ensures the paths being
checked use the prescribed infinite-time semantics.
\end{hbtheorem}

\subsection{Decision algorithms and their resource bounds}
\label{sec:p08-ltl-decision-bounds}

\begin{hbtheorem}{thm:p08-ltl-finite-graph-decision}{proved=annex:p08.ltl.finite-graph-decision}{A complete finite-graph emptiness algorithm}
Reachability followed by a two-pass depth-first decomposition
into strongly connected components decides automaton emptiness.
It accepts exactly the reachable cyclic components meeting \(F\).
The graph computation takes time linear in its explicit numbers
of vertices and edges. Consequently explicit LTL construction
and product checking take time \(2^{O(n)}\) and
\(\lVert K\rVert2^{O(n)}\), respectively, with exponential
graph storage permitted.
\end{hbtheorem}

\begin{hbtheorem}{thm:p08-ltl-polynomial-space}{proved=annex:p08.ltl.polynomial-space}{An explicit deterministic polynomial-space decision procedure}
LTL satisfiability and validity are decidable in deterministic
polynomial space in the formula length. Combined finite-state
model checking is decidable in space
\(O((n+\log\lVert K\rVert)^2)\), apart from the read-only
input. The procedure recomputes bounded reachability rather than
storing the exponential graph; this space claim does not require
that recomputation to retain the explicit-graph time bound.
\end{hbtheorem}

\subsection{The matching lower bounds}
\label{sec:p08-ltl-lower-bounds}

\begin{hbdefinition}{def:p08-ltl-configuration-encoding}{A finite formula describing polynomial-space configurations}
For this complexity statement atom names come from an unbounded
effectively named supply. Fix a deterministic polynomial-space
Turing-machine decider \(M\). On input \(w\), choose a
polynomially bounded width \(s\geq3\) large enough for its
tape, including two permanent boundary cells. Acceptance and
rejection have stationary self-transitions after halting.
Let \(\mathcal D_M\) be the finite cell alphabet consisting of a
boundary symbol, ordinary tape symbols and pairs of a machine
state with a tape symbol. Such a pair marks the unique head.

Use propositions \(C_0,\ldots,C_{s-1}\) for position within
a configuration and \(P_d\) for each \(d\in\mathcal D_M\).
For a nonempty finite list define the Boolean abbreviation
\[
 \operatorname{One}(B_1,\ldots,B_k)=
 (\bigvee_i B_i)\land
 \bigwedge_{i<j}\neg(B_i\land B_j).
\]
Let \(H\) be the disjunction of head-symbol propositions and
\(H_{\rm acc}\) that of accepting-head propositions. The
formula \(\Psi_{M,w}\) is the conjunction of these finite conditions:
\begin{enumerate}
\item \(C_0\) holds initially, the position propositions are
      always exactly one, and
      \(\mathsf G(C_i\to\mathsf XC_{i+1\bmod s})\)
      holds for each \(i\).
\item Cell propositions are always exactly one. The boundary
      proposition holds exactly at \(C_0\) and \(C_{s-1}\).
\item At each \(C_0\), exactly one of
      \(H,\mathsf XH,\ldots,\mathsf X^{s-1}H\) holds.
\item For each \(0\leq i<s\),
      \(\mathsf X^iP_{d_i}\) gives the known initial
      configuration: input \(w\), initial head/state, blank
      unused cells and the two boundary symbols.
\item Let \(f:\mathcal D_M^3\to\mathcal D_M\) be the local next-cell
      function of the machine. For each interior cell
      \(1\leq i\leq s-2\) and triple \((a,b,c)\), include
      \[
       \mathsf G\left(C_{i-1}\to
       \bigl((P_a\land\mathsf XP_b\land\mathsf X^2P_c)
                  \to\mathsf X^{s+1}P_{f(a,b,c)}\bigr)\right).
      \]
\item Finally include \(\mathsf F H_{\rm acc}\).
\end{enumerate}
The local function is specified completely in the proof. The
Boolean abbreviations and finite iterations of next are expanded,
so this is an ordinary formula of the registered language.
\end{hbdefinition}

\begin{hbtheorem}{thm:p08-ltl-pspace-hardness}{proved=annex:p08.ltl.pspace-hardness}{Satisfiability and validity are PSPACE-complete}
For each fixed decider \(M\), the formula \(\Psi_{M,w}\)
is constructible in polynomial time in \(\operatorname{len}(w)\), and is
satisfiable exactly when \(M\) accepts \(w\). Thus LTL
satisfiability is PSPACE-hard. With the upper bound above it
is PSPACE-complete, as is validity. This concerns unrestricted
Boolean negation; it does not assign the same complexity to
every positive or restricted temporal fragment.
\cite[Section 4, especially Lemma 4.4]{SistlaClarke1985}
\end{hbtheorem}

\begin{hbdefinition}{def:p08-ltl-block-translation}{Encoding arbitrary valuations in a small transition model}
List the input formula's atoms as \(p_0,\ldots,p_{k-1}\),
adding an unused atom if necessary so \(k\geq1\).
Construct states \((i,b)\), \(0\leq i<k,b\in\{0,1\}\),
and a fresh initial state \(w_*\). From \(w_*\) go to
either \((0,0)\) or \((0,1)\); from \((i,b)\) go to
either \((i+1\bmod k,0)\) or \((i+1\bmod k,1)\).
Use two fresh propositions: \(c\) holds exactly in phase zero,
and \(b_*\) holds exactly when the state's second coordinate
is one. The initial state carries neither.

Define \(T\) on formulas by preserving falsity, negation and
conjunction, and by
\[
\begin{aligned}
 T(p_i)&=\mathsf X^i b_*,\\
 T(\mathsf XA)&=\mathsf X^kT(A),\\
 T(A\mathbin{\mathsf U}B)
   &=(\neg c\lor T(A))\mathbin{\mathsf U}(c\land T(B)).
\end{aligned}
\]
At phase-zero positions the next \(k\) bits encode one valuation
of the original atoms. The until witness is forced to a phase-zero
position, and intermediate nonzero phases impose no left-operand
obligation.
\end{hbdefinition}

\begin{hbtheorem}{thm:p08-ltl-model-checking-complete}{proved=annex:p08.ltl.model-checking-complete}{Combined finite-state model checking is PSPACE-complete}
The constructed transition structure satisfies
\(\mathsf XT(A)\) on all paths from \(w_*\) exactly when
\(A\) is valid over the original atom valuations. Its size is
\(O(k)\) and the translated formula size is
\(O(k\operatorname{len}(A))\), where this instance of
\(\operatorname{len}\) counts written formula symbols.
This polynomial reduction proves PSPACE-hardness of universal
finite-state LTL checking; together with the proved upper bound
it gives PSPACE-completeness. The initial next operator skips
only the fresh branching start, not any encoded valuation.
\end{hbtheorem}

\paragraph{Reference use.}
For satisfiability, construct the formula's local types and pending
batches, then test for a reachable accepting positive cycle.
For validity, run that test on the syntactic negation: an
accepted lasso is a counterexample word. For a finite-state
specification, test the product with the automaton for its
negation; a returned lasso is a counterexample program path.
In each case the word or path uses the exact current-witness
until convention. An empty negation automaton proves validity,
and an empty counterexample product proves universal satisfaction.
The finite-type,
pending-batch and generalized-acceptance representations support
the same future-LTL satisfiability and verification tasks.
Strict tense, past-operator, determinization and branching-time
constructions require their respective languages and acceptance
conditions.
\index{linear temporal logic!completeness}\index{linear temporal logic!decision}
\index{B\"uchi automaton!eventuality obligations}\index{noncompactness!LTL}
